%-------------------------------------------------------------------------------
% §2 MATRICE DE CONFORMITÉ AUX EXIGENCES DES TDR
%-------------------------------------------------------------------------------
\cvsection{Matrice de conformité aux exigences des TDR}

\renewcommand{\arraystretch}{1.4}
\begin{longtable}{>{\bfseries\raggedright\arraybackslash}p{6.4cm} >{\raggedright\arraybackslash}p{10.3cm}}
  \toprule
  \textbf{Exigence TDR (Section IX)} & \textbf{Preuve} \\
  \midrule
  \endfirsthead

  \toprule
  \textbf{Exigence TDR (Section IX)} & \textbf{Preuve} \\
  \midrule
  \endhead

  \bottomrule
  \endlastfoot

  Master ou Ingéniorat en informatique / génie logiciel
    & Ingénieur (Bac+5 / niveau Master) en Réseaux et Télécommunications — UTT, France — voir §3 \\

  10 ans d'expérience conception / développement SI
    & Plus de 10 ans : Covage (2015) → présent — voir §5 \\

  Dont ≥ 5 ans secteur public
    & $\sim$7 ans : Covage (RIP / délégation de service public, 2015–2017) · Doctolib (téléservice public national, 2020–2022) · Nemidis / OBR (institution publique nationale, 2024–présent) \\

  Technologies de développement (PHP, JS, frameworks)
    & Ruby on Rails · PHP · React · Next.js · Node.js · TypeScript · Docker · CI/CD \\

  SGBD
    & PostgreSQL · MySQL · SQLite · Redis (NoSQL) \\

  Intégration API REST (≥ 3 ans)
    & REST · JSON · SFTP · Webhooks · Microservices — Climb (2023–2025), Doctolib (2020–2022), Joinly (2018–2019) \\

  2 audits d'applications (< 3 ans)
    & Audit NEMIDIS POS (2025) · Audit architecture des accès Abbove (2026) — voir §6 \\

  Supervision ≥ 2 projets SI d'envergure
    & OBR / MFE / ASYCUDA–EBMS (2025–2026) · Doctolib e-Prescription (2020–2022) · Climb Data Platform (2023–2025) — voir §4 \\

  ≥ 1 cahier des charges / architecture globale d'un SI
    & Architecture ASYCUDA–EBMS–POS-OBR · Architecture CMMS Nemidis — voir §7 \\

  Français et Anglais courants
    & Courant (professionnel et quotidien) \\

  Kirundi (avantage)
    & Langue maternelle \\

  Collaboration avec institutions publiques
    & OBR (Office Burundais des Recettes) — institution publique nationale \\

  \bottomrule
\end{longtable}
